\section{What a cost budget can guarantee}\label{sec:allocation}

The allocation results in this section are elementary accounting and
scheduling facts. They place useful limits on the performance claims one
can derive from a guarded enhancement. No publication priority is claimed
for them.

\begin{proposition}[Cost on the enhanced execution]\label{prop:ledger}
Let $N$ be the time charged to native work on one enhanced execution and
$E$ the time charged to optional enhancement work. Admit a new enhancement
only when $E<b+\eta N$, where $b,\eta\ge0$. If each admitted operation
adds at most $\delta$ to $E$, and every increment of $E$ belongs to such
an admitted operation, then at every completed operation
\begin{equation}\label{eq:ledger}
 E\le b+\eta N+\delta,\qquad
 N+E\le(1+\eta)N+b+\delta.
\end{equation}
If each cost has an enforced upper bound $r$ and admission instead
requires $E+r\le b+\eta N$, the final overshoot $\delta$ can be removed.
\end{proposition}
\begin{proof}
Immediately before the last admitted operation, the admission condition
holds. That operation adds at most $\delta$; later native work only
increases the right-hand side. If no operation was admitted, the result
is immediate. An advance reservation that includes the complete eventual
cost prevents the final overshoot.
\end{proof}

The charge must include discovery, model assembly, screening, optimization,
verification, and callback bookkeeping. A time limit on an auxiliary LP
does not bound surrounding Python or symbolic work. An implementation
without an enforced bound on the whole operation must record overshoot
rather than claim that~\eqref{eq:ledger} holds with its nominal LP limit.
If rejected entry or admission checks incur an unreserved cost $J$, that
cost must be added separately to the total. The proposition gives no
bound on $J$ without an additional premise limiting it.

The quantity $N$ refers to native work along the enhanced trajectory.
Let $N_0$ denote the work in an independent baseline run. Nothing in
Proposition~\ref{prop:ledger} connects $N$ and $N_0$. An abstract example
makes the missing premise clear. At the first decision a deterministic
search procedure has two legal choices. Its default ordering finds a
complete proof in one unit. Adding a redundant row changes the ordering
and leads to a proof after $M$ units. The row costs $b>0$ to obtain.
The enhancement budget is respected, but the runtime ratio is $M+b$,
which is unbounded as $M$ grows. This is a counterexample in an abstract
search model, not a constructed SCIP instance or a prediction of SCIP
behavior.

An independent, unchanged baseline provides a different guarantee.

\begin{proposition}[Protected baseline scheduling]\label{prop:portfolio}
In an ideal preemptible work model, run independent solvers $A$ and $H$
with separate states and fixed random tapes. Suppose by total work time
$t$, solver $A$ has received at least $(1-\eta)t-h$ work and $H$ at
least $\eta t-h$, for $0<\eta<1$ and $h\ge0$. If their standalone
completion requirements are $T_A,T_H$, the first answer arrives by
\begin{equation}\label{eq:portfolio}
 \min\left\{\frac{T_A+h}{1-\eta},\frac{T_H+h}{\eta}\right\}.
\end{equation}
Alternatively, run $H$ for at most $b$ charged work, then, if necessary,
start unchanged $A$ from its initial state. The total work is at most
$b+T_A$, provided startup is charged.
\end{proposition}
\begin{proof}
At either time in~\eqref{eq:portfolio}, the corresponding solver has
received its required work. For serial rescue, either $H$ finishes within
$b$, or its cost at most $b$ is followed by exactly the independent
baseline execution.
\end{proof}

The assumptions exclude interference from shared memory exhaustion,
cache contention, uncharged startup, and wall-clock-dependent decisions.
Injecting a new incumbent or rows into the baseline can change its work
requirement and needs a separate argument. This report does not deploy
a solver portfolio or claim the scheduling bound for its experiments.

There is also no universally correct first action from an arbitrary
reward history. Suppose two actions each cost one unit and two possible
environments have identical observable histories. Only the first action
reveals a proof in one environment; only the second does so in the other.
A deterministic selector fails to choose correctly first in one
environment. A randomized first choice succeeds with probability at most
one half in at least one environment. Structural information or a
statistical assumption can distinguish the cases; history alone cannot.
This observation does not rule out learning or competitive scheduling.

These limits suggest a small interface for optional solver actions.
Each action declares its model, node domain, relaxation, cutoff, and cost
allowance. It returns a verified inference, a verified upper bound on
benefit in a stated mathematical measure, or an unresolved status.
A selector may rank unresolved actions heuristically, but must not turn
a ranking into a validity claim. The ledger charges selection as well as
optimization. An OBBT width reduction, a support-cut violation, and a
decomposition lower-bound improvement are different measures; converting
them to expected seconds saved requires a separately validated model.
The present implementation and comparison concern the OBBT case.
